\section{Related work}

\paragraph{Streaming speech translation.}
Streaming SST balances quality against latency while decoding partial speech.
Prefix-to-prefix policies make this trade-off explicit \citep{ma2019stacl},
attention-based policies refine read and write decisions
\citep{papi2023alignatt,papi2024streamatt}, and SimulEval standardizes
evaluation \citep{ma2020simuleval}. Speech-language models broaden direct SST
architectures \citep{seamless2023,zhang2024streamspeech,xu2025qwen3omni}.
AutoTerm-SST keeps an existing streaming model and adds terminology management
and an inspection layer around it.

\paragraph{Terminology-aware translation.}
Terminology-aware methods take the relevant glossary as given before
decoding. Rare terms and named entities remain a weakness of speech translation
\citep{gaido2021moby}, and prior work applies user-provided dictionaries
through constrained decoding or training-time exposure
\citep{hokamp2017lexically,dinu2019training}. Retrieval-augmented generation
offers a non-parametric alternative \citep{lewis2020retrieval,khandelwal2021nearest}, which RASST
applies to streaming LLM decoding \citep{luo2026rasst}. AutoTerm-SST removes the
pre-session glossary by selecting topic slices online, and like interactive
translation tools \citep{dandekar2024translation} it makes these decisions
visible, on top of a continuous-batching serving engine
\citep{kwon2023efficient,zheng2024sglang}.
